\scanpage{77}
\mspage{73}
partie $\mathbb{\Sigma}_{\underline{E}}$ (ou $\mathbb{\Sigma}_{E}$) de $\unsure{\mathcal{P}|\mathrm{P}}_{E}$, stable par $\bigcap_{i}$ \unsure{quelconque|quelconques}, ayant comme \unsure{pl. grand|pl. grand.} \unsure{élémt|élément} $E$, comme plus petit élément le \unsure{sous-ensemble} $E_{0}$. \deleted{\unsure{Il}}\unsure{Elle|Il} a des $\operatorname{Inf}$ quelc. (les $\bigcap_{i}$) et $\operatorname{Sup}$ quelc. (en général $\deleted{\unsure{\not\leq|\neq}}\deleted{\unsure{les}}\ne\bigcup_{i}$), et une anti-involution
\[
C\deleted{\unsure{A}}\unsure{S|A}\overset{\mathrm{def}}{=}\operatorname{cosupp}_{\deleted{\unsure{E}}}^{\circ}A
\qquad (A\in\mathbb{\Sigma}_{E}\subset\mathcal{P}(E)).
\]

Application canonique
\[
x\longmapsto\sigma(x)=\operatorname{supp}_{E}^{\circ}(\{x\}):E\xrightarrow{\sigma}\mathbb{\Sigma}_{E}
\]
% margin note written sideways
\margin{On a $x\mathrel{|\unsure{\circ|0}|}y\Longleftrightarrow\sigma x\mathrel{|\circ|}\sigma y$ où pour $S,T\in\mathbb{\Sigma}_{\unsure{E}}$, plus \unsure{génér.} pour $A,B\in\mathcal{P}(E)$, on pose $A\mathrel{|\circ|}B\Longleftrightarrow x\mathrel{|\circ|}y\quad\forall x\in A,y\in B$. \unsure{Or.|On.} \deleted{Si $S,T\in\mathbb{\Sigma}_{E}$, \unsure{on}} $A\mathrel{|\circ|}B\Longleftrightarrow A\subset\operatorname{cosupp}^{\circ}B\Longleftrightarrow B\subset\operatorname{cosupp}^{\circ}A\Longleftrightarrow\operatorname{supp}^{\circ}(A)\subset\operatorname{cosupp}^{\circ}B\Longleftrightarrow\operatorname{supp}^{\circ}B\subset\operatorname{cosupp}^{\circ}A$. Dans $\mathbb{\Sigma}_{E}$, la relation $\mathrel{|\circ|}$ \deleted{\unsure{de}} résulte donc de la structure $\leq$ et $C$.}

\uline{Disjoncteur} \uline{trivial} : quand
\[
x\mathrel{|\circ|}y\Longleftrightarrow x\ne y
\]
donc cela correspond \unsure{:|à} $A=\emptyset$, structure de graphe \unsure{discrète|discrèt}. On a alors
\[
\mathbb{\Sigma}_{E}=\mathcal{P}(E).
\]

\uline{NB} Pour toute partie $A$ de $E$, on a
\[
\begin{aligned}
\operatorname{cosupp}_{E}^{\circ}(A)&=\operatorname{cosupp}_{E}^{\circ}(A\cup E_{0})\supset E_{0},\qquad\text{d'où}\\
\deleted{\unsure{\mathrm{c}}}\operatorname{Supp}_{E}^{\circ}(\unsure{M|A})&=\operatorname{supp}_{E}^{\circ}(A\cup E_{0})\unsure{>|\supset} E_{0},
\end{aligned}
\]
et par suite aussi
\[
\begin{aligned}
\operatorname{cosupp}_{E}^{\circ}(A)&=\operatorname{cosupp}^{\circ}(A^{\unsure{\ast|\prime}})\\
\operatorname{supp}_{E}^{\circ}(M)&=\operatorname{supp}^{\circ}(A^{\ast})
\end{aligned}
\]
où
\[
A^{\ast}=A\setminus E_{0}\cap A.
\]
On en conclut que l'on a un iso. (respectant structure d'ordre et $C$)
\[
\begin{aligned}
\mathbb{\Sigma}_{E}\mathrel{\overset{\sim}{\longleftrightarrow}}\mathbb{\Sigma}_{E^{\ast}}\qquad &\deleted{\unsure{A}}S\longmapsto S\setminus E_{0}\\
&S^{\ast}\mathbin{\unsure{\sqcup|\cup}}E_{0}\mathrel{\longleftarrow\mkern-9mu\vert} S^{\ast}.
\end{aligned}
\]
% archive sheet number: 76
